\documentclass[12pt,a4paper]{ctexart}

% === 页面设置 ===
\usepackage[top=2.5cm,bottom=2.5cm,left=2.5cm,right=2.5cm]{geometry}

% === 数学和物理 ===
\usepackage{amsmath,amssymb,amsthm}
\usepackage{physics}

% === 图形和表格 ===
\usepackage{graphicx}
\usepackage{float}
\usepackage{booktabs}
\usepackage{array}
\usepackage{caption}
\usepackage{subcaption}

% === 代码 ===
\usepackage{listings}
\usepackage{xcolor}
\usepackage{upquote}

\definecolor{codebg}{rgb}{0.95,0.95,0.95}
\definecolor{codeframe}{rgb}{0.7,0.7,0.7}
\definecolor{commentcolor}{rgb}{0.3,0.5,0.3}
\definecolor{stringcolor}{rgb}{0.7,0.3,0.1}
\definecolor{keywordcolor}{rgb}{0.2,0.2,0.7}

\lstset{
    basicstyle=\ttfamily\small,
    backgroundcolor=\color{codebg},
    frame=single,
    rulecolor=\color{codeframe},
    breaklines=true,
    numbers=left,
    numberstyle=\tiny\color{gray},
    xleftmargin=2em,
    framexleftmargin=1.5em,
    commentstyle=\color{commentcolor},
    stringstyle=\color{stringcolor},
    keywordstyle=\color{keywordcolor}\bfseries,
    showstringspaces=false,
    tabsize=4,
    captionpos=b,
}

% === 引用 ===
\usepackage{hyperref}
\usepackage[numbers,sort&compress]{natbib}

% === 其他 ===
\usepackage{enumitem}
\usepackage{fancyhdr}
\usepackage{multirow}

\hypersetup{
    colorlinks=true,
    linkcolor=blue,
    citecolor=blue,
    urlcolor=blue,
}

% === 标题信息 ===
\title{\heiti PyQUDA \\[3pt] \large QUDA的Python/Cython接口使用手册}
\author{张鑫}
\date{\today}

\begin{document}

\maketitle

\begin{abstract}
\noindent
PyQUDA是QUDA（用于格点QCD的GPU加速库）的Python/Cython封装。它结合了Python的易用性和QUDA的GPU高性能计算能力，支持多GPU分布式计算（通过MPI for Python），并提供完整的格点QCD计算功能：夸克传播子求解（Wilson/Clover/HISQ费米子作用量，含多重网格支持）、混合蒙特卡洛（HMC）、规范固定、规范涂抹、费米子涂抹、梯度流以及10余种格点数据格式的文件I/O。本手册详细介绍PyQUDA的安装配置、三层架构、核心功能、示例演示以及完整的API接口清单。
\end{abstract}

\tableofcontents
\newpage

%=============================================================================
\section{项目概述}
%=============================================================================

PyQUDA是一个用Cython编写的QUDA Python封装库，由CLQCD合作组开发和维护。该项目旨在利用Python中的优化线性代数库（如CuPy、PyTorch、DPNP）与QUDA的GPU计算能力相结合，实现高性能的格点QCD研究运算。

\subsection{主要特性}
\begin{itemize}[leftmargin=*]
    \item \textbf{多GPU支持}：通过mpi4py实现MPI分布式计算，支持4D进程网格拓扑
    \item \textbf{多后端GPU加速}：支持CuPy（NVIDIA CUDA/AMD HIP）、PyTorch（CUDA/HIP/SYCL）、DPNP（Intel GPU）及numpy（CPU回退）
    \item \textbf{丰富的文件I/O}：支持Chroma、MILC、KYU、XQCD、NERSC、OpenQCD、ILDG、NPY、HDF5等10余种格式
    \item \textbf{多种夸克作用量}：Wilson、Clover（各向同性/各向异性）、HISQ费米子
    \item \textbf{多重网格求解器}：支持Wilson和Clover费米子的多重网格加速
    \item \textbf{HMC模拟}：纯规范场、1/2味Clover费米子、N味HISQ费米子
    \item \textbf{规范操作}：规范固定、规范涂抹（APE/Stout/HYP）、梯度流（Wilson/Symanzik）
    \item \textbf{费米子涂抹}：高斯/Wuppertal涂抹
\end{itemize}

\subsection{三层架构}

PyQUDA采用清晰的三层架构设计：

\begin{figure}[H]
\centering
\fbox{\parbox{0.9\textwidth}{\centering
\vspace{0.3cm}
\textbf{Layer 2: pyquda\_utils} — 高级工具脚本 \\
\hspace{1em}core, source, gamma, phase, wilson\_loop, checksum, convert, fft \\[5pt]
\textbf{Layer 2: pyquda\_io} — 文件格式I/O \\
\hspace{1em}chroma, milc, kyu, xqcd, nersc, openqcd, ildg, npy, lime \\[5pt]
\textbf{Layer 2: pyquda\_plugins} — C/CUDA插件系统 \\
\hspace{1em}pycontract (夸克收缩), pygwu (overlap费米子) \\[8pt]
\rule{\linewidth}{0.5pt} \\[3pt]
\textbf{Layer 1: pyquda\_utils / pyquda\_io / pyquda\_plugins} \\
\rule{\linewidth}{0.5pt} \\[3pt]
\textbf{Layer 0: pyquda\_core} — Cython绑定（libquda） \\
\hspace{1em}pyquda (QUDA API), pyquda\_comm (MPI通信+数组后端)
\vspace{0.3cm}}}
\caption{PyQUDA 三层架构}
\label{fig:architecture}
\end{figure}

\subsection{PyPI发布}

两个独立的PyPI包：
\begin{itemize}[leftmargin=*]
    \item \textbf{PyQUDA}（核心层）：依赖mpi4py、numpy，提供pyquda和pyquda\_comm模块
    \item \textbf{PyQUDA-Utils}（工具层）：依赖PyQUDA，提供pyquda\_utils、pyquda\_io、pyquda\_plugins模块
\end{itemize}

%=============================================================================
\section{安装}
%=============================================================================

\subsection{前置要求}

\begin{table}[H]
\centering
\caption{安装前置要求}
\label{tab:prerequisites}
\begin{tabular}{@{}ll@{}}
\toprule
\textbf{组件} & \textbf{说明} \\
\midrule
QUDA & 编译安装的QUDA库（develop分支，\texttt{\$QUDA\_PATH}指向build/usqcd） \\
CUDA/ROCm & GPU计算工具链 \\
MPI & 支持GPU直连的MPI实现（如OpenMPI） \\
Python & $\ge$ 3.8 \\
pip依赖 & numpy, mpi4py, cupy (或torch/dpnp) \\
编译依赖 & Cython, setuptools, wheel \\
\bottomrule
\end{tabular}
\end{table}

\subsection{从PyPI安装（推荐）}

\begin{lstlisting}[caption={从PyPI安装}, label={lst:pypi_install}]
# 确保QUDA_PATH指向QUDA安装目录
export QUDA_PATH=/path/to/quda/build/usqcd

# 安装核心包和工具包
python3 -m pip install pyquda pyquda-utils
\end{lstlisting}

\subsection{从源码安装}

\begin{lstlisting}[caption={从源码编译安装}, label={lst:source_install}]
export QUDA_PATH=/path/to/quda/build/usqcd

# 克隆仓库（含子模块）
git clone --recursive https://github.com/CLQCD/PyQUDA.git
cd PyQUDA

# 安装核心Cython绑定
cd pyquda_core
python3 -m pip install .
cd ..

# 安装工具层
python3 -m pip install .
\end{lstlisting}

\subsection{开发模式安装（原地构建）}

\begin{lstlisting}[caption={原地构建（开发模式）}, label={lst:dev_install}]
export QUDA_PATH=/path/to/quda/build/usqcd
git clone --recursive https://github.com/CLQCD/PyQUDA.git
cd PyQUDA

# 原地构建Cython扩展
cd pyquda_core
python3 setup.py build_ext --inplace
cd ..

# 创建符号链接
python3 setup.py egg_info
ln -s pyquda_core/pyquda_comm ./
ln -s pyquda_core/pyquda ./
\end{lstlisting}

注意：从其他目录运行脚本时需将PyQUDA根目录添加到\texttt{sys.path}。

%=============================================================================
\section{初始化与网格拓扑}
%=============================================================================

\subsection{基本初始化}

\begin{lstlisting}[caption={PyQUDA初始化}, label={lst:init}]
from pyquda import init
from pyquda_utils import core

# 方法1：简单初始化
core.init(resource_path=".cache/quda")

# 方法2：使用pyquda.init
from pyquda_comm import initGrid, initDevice
from pyquda import initQUDA

latt_info = core.LatticeInfo([24, 24, 24, 72])  # [Lx, Ly, Lz, Lt]
initGrid(latt_info.grid_size)                    # 4D进程网格
initDevice(0)                                    # 选择设备
initQUDA()                                       # 初始化QUDA
\end{lstlisting}

\subsection{LatticeInfo}

\begin{lstlisting}[caption={格点信息对象}, label={lst:lattice_info}]
from pyquda_utils.core import LatticeInfo

# 创建格点信息
latt_info = LatticeInfo(
    [24, 24, 24, 72],          # [Lx, Ly, Lz, Lt]
    grid_size=[1, 1, 1, 4],    # 4D进程网格
    t_boundary=-1,              # 时间边界条件 (-1: 反周期)
    anisotropy=1.0,             # 各向异性系数
)

# 访问属性
print(latt_info.Gs)   # 空间点总数
print(latt_info.Gt)   # 时间点总数
print(latt_info.volume)  # 总格点数
\end{lstlisting}

\subsection{进程网格映射类型}

支持以下网格映射策略：
\begin{itemize}[leftmargin=*]
    \item \texttt{default}：默认映射
    \item \texttt{t\_fastest}：时间维度最快变化
    \item \texttt{x\_fastest}：空间维度最快变化
    \item \texttt{minimize}：最小化通信面
    \item \texttt{shared}：共享内存映射
    \item \texttt{dist\_graph}：分布式图映射
\end{itemize}

%=============================================================================
\section{格点场类型系统}
%=============================================================================

PyQUDA提供了分层的格点场类型系统。

\subsection{场类型层次}

\begin{lstlisting}[caption={格点场类型}, label={lst:field_types}]
from pyquda_comm.field import (
    LatticeInfo,          # 格点几何元数据
    LatticeGauge,         # 规范场 U_mu(x): [Nc, Nc, Nd, Lx, Ly, Lz, Lt]
    LatticeFermion,       # 费米子场: [Ns, Nc, Lx, Ly, Lz, Lt]
    LatticePropagator,    # 传播子: [Ns, Nc, Ns, Nc, Lx, Ly, Lz, Lt]
    LatticeStaggeredFermion,  # 交错费米子
    LatticeMom,           # 动量场 (HMC)
    LatticeClover,        # Clover项
    LatticeRotation,      # 旋转矩阵 (规范固定)
    MultiLatticeFermion,  # 多源费米子
)
\end{lstlisting}

\subsection{场的基本操作}

\begin{lstlisting}[caption={场的创建和操作}, label={lst:field_ops}]
# 创建规范场
gauge = LatticeGauge(latt_info)

# 加载规范组态
gauge.loadNPYGauge("path/to/config.npy")
gauge.loadChromaGauge("path/to/config.lime")

# 保存规范场
gauge.saveNPYGauge("output.npy")

# 偶-奇预处理
gauge_even = gauge.even
gauge_odd = gauge.odd

# 规范化检查
checksum = gauge.checksum()
\end{lstlisting}

\subsection{场的类型（pyquda.field扩展）}

\texttt{pyquda/field.py}在\texttt{pyquda\_comm}的基础场上添加了高级方法：

\begin{lstlisting}[caption={LatticeGauge 高级方法}]
class LatticeGauge:
    # HDF5 I/O
    def loadHDF5Gauge(self, path): ...
    def saveHDF5Gauge(self, path): ...

    # 规范场操作（通过gauge_dirac属性）
    @property
    def gauge_dirac(self) -> GaugeDirac: ...

    # 涂抹
    def smearAPE(self, coeff, steps): ...
    def smearStout(self, coeff, steps): ...
    def smearHYP(self, coeff, steps): ...

    # 梯度流
    def wflow(self, steps, step_size): ...
    def symanzikFlow(self, steps, step_size): ...

    # 规范固定
    def gaugeFixingOVR(self, tolerance): ...
    def gaugeFixingFFT(self): ...

    # 物理观测量
    def plaquette(self): ...
    def polyakov(self): ...
    def topologicalCharge(self): ...
\end{lstlisting}

%=============================================================================
\section{Dirac算子}
%=============================================================================

\subsection{算子层次结构}

\begin{figure}[H]
\centering
\fbox{\parbox{0.9\textwidth}{\centering
\vspace{0.3cm}
\textbf{Dirac}（抽象基类） \\
$\downarrow$ \\
\textbf{FermionDirac} — Wilson/Clover费米子 + 多重网格 \\
$\downarrow$ \\
WilsonDirac / CloverWilsonDirac \\[5pt]
\textbf{StaggeredFermionDirac} — 交错费米子 \\
$\downarrow$ \\
StaggeredDirac / HISQDirac \\[5pt]
\textbf{GaugeDirac} — 规范场操作（涂抹、流动、固定）
\vspace{0.3cm}}}
\caption{Dirac算子继承层次}
\label{fig:dirac}
\end{figure}

\subsection{创建和使用Dirac算子}

\begin{lstlisting}[caption={Dirac算子使用示例}, label={lst:dirac}]
from pyquda_utils.core import getWilson, getClover, getHISQ

# Wilson费米子
wilson_dirac = getWilson(latt_info, gauge, mass=-0.2770)

# Clover费米子（自动创建Clover项）
clover_dirac = getClover(
    latt_info, gauge, mass=-0.2770,
    clover_coeff=1.0,            # clover系数
    multigrid=True,              # 启用多重网格
)

# HISQ费米子
hisq_dirac = getHISQ(latt_info, gauge, mass=0.05)
\end{lstlisting}

\subsection{传播子求解}

\begin{lstlisting}[caption={夸克传播子求解}, label={lst:invert}]
from pyquda_utils.core import invert
from pyquda_utils.source import wall

# 墙源传播子
prop = invert(
    clover_dirac,
    wall([24, 24, 24, 72], latt_info),  # 源矢量
)

# 多源传播子
props = invert(
    clover_dirac,
    [wall([24, 24, 24, 72], latt_info) for _ in range(4)],  # 4源
)

# 点源传播子
from pyquda_utils.source import point
prop = invert(clover_dirac, point(latt_info, [0, 0, 0, 0]))
\end{lstlisting}

\subsection{Dirac算子方法}

\begin{lstlisting}[caption={FermionDirac 关键方法}]
class FermionDirac:
    def invert(self, source, mass=None, tol=1e-12, maxiter=10000): ...
    def invertRestart(self, source, mass=None, ...): ...
    def invertPC(self, source, mass=None, ...): ...       # 偶-奇预处理
    def invertCloverPC(self, source, mass=None, ...): ...
    def invertMultiSrc(self, sources, mass=None, ...): ...  # 多源同时求解
    def invertMultiSrcRestart(self, sources, mass=None, ...): ...
    def mat(self, vector): ...            # 矩阵乘法
    def matDagMat(self, vector): ...      # D^dag * D
    def dslash(self, vector): ...         # Dslash算子
    def performance(self): ...            # 性能报告
\end{lstlisting}

\subsection{精度控制}

全局精度预设：
\begin{itemize}[leftmargin=*]
    \item \texttt{none}：不配置精度
    \item \texttt{invert}：配置反演精度（double/half/single混合）
    \item \texttt{multishift}：多shift求解精度
    \item \texttt{multigrid}：多重网格精度
\end{itemize}

全局重构预设（减少内存占用）：
\begin{itemize}[leftmargin=*]
    \item \texttt{none}：不重构
    \item \texttt{wilson}：Wilson费米子12/8重构
    \item \texttt{staggered}：交错费米子13/9重构
\end{itemize}

%=============================================================================
\section{HMC（混合蒙特卡洛）}
%=============================================================================

\subsection{HMC使用示例}

\begin{lstlisting}[caption={淬火HMC示例}, label={lst:hmc_example}]
from math import exp
from pyquda.action import GaugeAction
from pyquda.action.abstract import LoopParam
from pyquda.hmc import HMC, Integrator
from pyquda_utils import core
from pyquda_utils.core import X, Y, Z, T

# 初始化
core.init(resource_path=".cache/quda")
latt_info = core.LatticeInfo([16, 16, 16, 32])

# 定义规范作用量
u_0 = 0.855453; beta = 6.20
loop_param = LoopParam(
    [[X, Y, -X, -Y], [X, Z, -X, -Z], [Y, Z, -Y, -Z],
     [X, T, -X, -T], [Y, T, -Y, -T], [Z, T, -Z, -T]],
    [1.0, 1.0, 1.0, 1.0, 1.0, 1.0],
)

# 自定义积分器
class O2Nf1Ng0V(Integrator):
    def integrate(self, updateGauge, updateMom, t):
        dt = t / self.n_steps
        for _ in range(self.n_steps):
            updateMom(dt / 2)
            updateGauge(dt)
            updateMom(dt / 2)

# 创建HMC
hmc = HMC(latt_info,
    [GaugeAction(latt_info, loop_param, beta)],
    O2Nf1Ng0V(n_steps=100))

# 运行HMC
gauge = core.LatticeGauge(latt_info)
hmc.initialize(10086, gauge)

for i in range(2000):
    hmc.gaussMom()
    energy_old = hmc.momAction() + hmc.gaugeAction()
    hmc.integrate(1.0, 2e-14)
    energy = hmc.momAction() + hmc.gaugeAction()
    accept = hmc.accept(energy - energy_old)
    if accept or i < 500:  # 热化500步
        hmc.saveGauge(gauge)
    else:
        hmc.loadGauge(gauge)
\end{lstlisting}

\subsection{内置积分器}

\begin{table}[H]
\centering
\caption{内置辛积分器}
\label{tab:integrators}
\begin{tabular}{@{}ll@{}}
\toprule
\textbf{积分器} & \textbf{说明} \\
\midrule
\texttt{INT\_3G1F} & 3规范力/1费米子力 \\
\texttt{O2Nf1Ng0V} & 2阶辛，1费米子力，0规范力，速度Verlet \\
\texttt{O2Nf1Ng0P} & 同上，位置Verlet \\
\texttt{O2Nf2Ng0V} & 2阶辛，2费米子力 \\
\texttt{O2Nf2Ng0P} & 同上，位置Verlet \\
\texttt{O4Nf5Ng0V} & 4阶辛，5费米子力 \\
\texttt{O4Nf5Ng0P} & 同上，位置Verlet \\
\bottomrule
\end{tabular}
\end{table}

\subsection{HMC 类方法}

\begin{lstlisting}[caption={HMC 完整接口}]
class HMC:
    def __init__(self, latt_info, monomials, integrator): ...

    def initialize(self, seed, gauge): ...
    def gaussMom(self): ...
    def samplePhi(self): ...
    def integrate(self, t, tol): ...
    def accept(self, dE) -> bool: ...

    # 作用量计算
    def gaugeAction(self) -> float: ...
    def fermionAction(self) -> float: ...
    def momAction(self) -> float: ...

    # 观测量
    def plaquette(self) -> float: ...
    def polyakovLoop(self) -> float: ...
    def energy(self) -> float: ...
    def qcharge(self) -> float: ...
    def qchargeDensity(self) -> float: ...

    # 规范场管理
    def saveGauge(self, gauge): ...
    def loadGauge(self, gauge): ...
\end{lstlisting}

\subsection{费米子作用量}

\begin{lstlisting}[caption={费米子作用量}]
# Clover-Wilson费米子（1/2味）
from pyquda.action.clover_wilson import CloverWilsonAction
clover_action = CloverWilsonAction(latt_info, gauge, mass, clover_coeff, n_flavor=2)

# HISQ费米子（N味）
from pyquda.action.hisq import HISQAction
hisq_action = HISQAction(latt_info, gauge, mass, n_flavor=4)

# 多重HISQ（共享fat/long链接）
from pyquda.action.hisq import MultiHISQAction
multi_hisq = MultiHISQAction(latt_info, gauge, [mass1, mass2], n_flavor=2)
\end{lstlisting}

%=============================================================================
\section{规范操作}
%=============================================================================

\subsection{规范涂抹}

通过\texttt{GaugeDirac}接口提供AP、Stout和HYP涂抹。

\begin{lstlisting}[caption={规范涂抹}]
# 获取GaugeDirac
gd = gauge.gauge_dirac

# APE涂抹
gauge.smearAPE(coeff=0.5, steps=10)
gauge.smearStout(coeff=0.1, steps=10)
gauge.smearHYP(coeff=0.5, steps=10)
\end{lstlisting}

\subsection{梯度流}

\begin{lstlisting}[caption={梯度流}]
# Wilson流
gauge.wflow(steps=100, step_size=0.01)

# Symanzik流
gauge.symanzikFlow(steps=100, step_size=0.01)

# 流标度确定 (t0, w0)
from pyquda.dirac.gauge import GaugeDirac
gd = GaugeDirac(latt_info)
gd.wflow_scale_determination(gauge, steps=100, step_size=0.01)
\end{lstlisting}

\subsection{规范固定}

\begin{lstlisting}[caption={规范固定}]
# 过松弛方法 (OVR)
gauge.gaugeFixingOVR(tolerance=1e-12)

# 快速傅里叶变换方法 (FFT)
gauge.gaugeFixingFFT()
\end{lstlisting}

\subsection{物理观测量}

\begin{lstlisting}[caption={规范观测量}]
# Plaquette
plaq = gauge.plaquette()

# Polyakov圈
poly = gauge.polyakov()

# 拓扑荷
Q = gauge.topologicalCharge()

# 规范作用量密度
E = gauge.energy()
\end{lstlisting}

%=============================================================================
\section{文件I/O}
%=============================================================================

PyQUDA支持10余种格点数据格式的读写。

\subsection{支持的格式}

\begin{table}[H]
\centering
\caption{文件I/O格式支持矩阵}
\label{tab:io_formats}
\begin{tabular}{@{}lcccc@{}}
\toprule
\textbf{格式} & \textbf{规范读} & \textbf{规范写} & \textbf{传播子读} & \textbf{传播子写} \\
\midrule
Chroma (SciDAC/LIME) & $\checkmark$ & — & $\checkmark$ & — \\
MILC & $\checkmark$ & $\checkmark$ & $\checkmark$ & — \\
KYU & $\checkmark$ & $\checkmark$ & $\checkmark$ & $\checkmark$ \\
XQCD & — & — & $\checkmark$ & $\checkmark$ \\
NERSC & $\checkmark$ & $\checkmark$ & — & — \\
OpenQCD & $\checkmark$ & $\checkmark$ & — & — \\
ILDG & $\checkmark$ & — & — & — \\
NPY (numpy) & $\checkmark$ & $\checkmark$ & $\checkmark$ & $\checkmark$ \\
HDF5 (h5py) & $\checkmark$ & $\checkmark$ & $\checkmark$ & $\checkmark$ \\
General (自定义) & $\checkmark$ & $\checkmark$ & — & — \\
\bottomrule
\end{tabular}
\end{table}

\subsection{加载/保存规范组态}

\begin{lstlisting}[caption={规范组态I/O}]
from pyquda_io import chroma, milc, npy

# Chroma格式
gauge = chroma.readChromaGauge("config.lime", latt_info)

# MILC格式
gauge = milc.readMILCGauge("config.milc", latt_info)

# KYU格式
from pyquda_io import kyu
kyu.writeKYUGauge("config.kyu", gauge)

# NPY格式
npy.writeNPYGauge("config.npy", gauge)
gauge = npy.readNPYGauge("config.npy", latt_info)

# HDF5格式
gauge.saveHDF5Gauge("config.h5")
gauge.loadHDF5Gauge("config.h5")
\end{lstlisting}

\subsection{校验和}

\begin{lstlisting}[caption={校验和验证}]
from pyquda_utils.checksum import checksum

# SciDAC风格校验和（CRC32 + rank-based XOR）
sum_a, sum_b = gauge.checksum()
\end{lstlisting}

%=============================================================================
\section{Gamma矩阵系统}
%=============================================================================

\begin{lstlisting}[caption={Gamma矩阵运算}]
from pyquda.gamma import Gamma, Projector

# 基本Gamma矩阵
g5 = Gamma(15)     # gamma_5
g0 = Gamma(0)      # gamma_t (DeGrand-Rossi基)
g1 = Gamma(1)      # gamma_x
g2 = Gamma(2)      # gamma_y
g3 = Gamma(3)      # gamma_z

# Gamma矩阵运算
g_plus = (g0 + g3) / 2    # (1+gamma_3)/2 投影算符
g5_g0 = g5 @ g0           # gamma_5 * gamma_t
g5g0_T = g5_g0.T          # 转置
g5g0_H = g5_g0.H          # 厄米共轭

# 自定义投影算符
proj = Projector([
    0.5, 0, 0, 0.5,    # 第一行
    0, 0, 0, 0,
    0, 0, 0, 0,
    0.5, 0, 0, 0.5,    # 最后一行
])  # (1+gamma_3)/2
\end{lstlisting}

\subsection{Gamma × 传播子}

\texttt{pyquda\_utils/gamma.py}通过monkey-patching为\texttt{LatticeFermion}和\texttt{LatticePropagator}添加矩阵乘法：

\begin{lstlisting}[caption={Gamma × 传播子操作}]
from pyquda_utils.gamma import *  # 自动添加 __matmul__ 到场类型

# gamma_5 @ 传播子
result = Gamma(15) @ prop

# 投影算符 @ 费米子场
projected = Projector(...) @ fermion
\end{lstlisting}

%=============================================================================
\section{源构造}
%=============================================================================

\texttt{pyquda\_utils/source.py}提供了多种夸克源构造方法：

\begin{lstlisting}[caption={源构造方法}]
from pyquda_utils.source import (
    point, wall, volume, momentum, colorvector, source, source12, source3
)

# 点源
src = point(latt_info, position=[0, 0, 0, 0])

# 墙源（时间片0上的所有空间点）
src = wall([24, 24, 24, 72], latt_info, t=0)

# 动量源（带Fourier相位）
src = momentum(latt_info, mom=[0, 0, 1, 0])

# 体积源
src = volume([24, 24, 24, 72], latt_info)

# 颜色矢量源
src = colorvector(latt_info, color=0, spin=0)

# 一般源（12分量: 自旋*颜色）
src = source12(latt_info)  # 无涂抹
src = source12(latt_info, source_type="gaussian")  # 高斯涂抹
\end{lstlisting}

\subsection{高斯/Wuppertal涂抹}

\begin{lstlisting}[caption={高斯涂抹}]
# 高斯涂抹源
from pyquda.dirac.gauge import GaugeDirac
gd = GaugeDirac(latt_info)
gd.gaussianSmearSource(source, gauge, n_steps=50, sigma=4.0)
\end{lstlisting}

%=============================================================================
\section{插件系统}
%=============================================================================

PyQUDA提供了一个自动化的C/CUDA插件系统，从C头文件自动生成Cython封装。

\subsection{内置插件}

\begin{itemize}[leftmargin=*]
    \item \textbf{pycontract}：GPU加速的介子/重子收缩函数
    \item \textbf{pygwu}：Grid-Wilson-Overlap接口
\end{itemize}

\subsection{pycontract — 夸克收缩}

\begin{lstlisting}[caption={pycontract 收缩函数}]
from pyquda_plugins.pycontract import (
    mesonTwoPoint, mesonAllSinkTwoPoint,
    baryonTwoPoint, baryonDiquark,
    baryonSequentialTwoPoint,
)

# 介子两点函数
C2pt = mesonTwoPoint(prop, prop, gamma_list)

# 重子两点函数
C2pt_baryon = baryonTwoPoint(prop, prop, gamma_list)
\end{lstlisting}

\subsection{pygwu — Overlap费米子}

\begin{lstlisting}[caption={pygwu Overlap接口}]
from pyquda_plugins.pygwu import Overlap

ov = Overlap(latt_info, gauge, mass=-0.2770, kappa=0.15)
ov.buildWilsonKernel()
ov.invert(source, mass=0.01, tol=1e-12)
\end{lstlisting}

%=============================================================================
\section{示例演示}
%=============================================================================

\subsection{示例1：淬火HMC}

\texttt{examples/1\_Quenched\_HMC.py} — 纯规范场HMC模拟，在$16^3\times32$格点上，$\beta=6.20$，使用Wilson规范作用量和O2Nf1Ng0V积分器。

\begin{lstlisting}[caption={运行淬火HMC}]
cd PyQUDA/examples
python 1_Quenched_HMC.py
# 输出: 每5条轨迹保存NPY格式规范组态到 ./cfg/
\end{lstlisting}

\subsection{示例2：夸克传播子与π介子两点函数}

\texttt{examples/2\_Quark\_Propagator.py} — 读取Chroma规范组态，施加Stout涂抹，使用多重网格Clover-Wilson Dirac求解器，从墙源传播子计算π介子两点函数。

\subsection{示例3：π介子和质子两点函数}

\texttt{examples/3\_Pion\_Proton\_2pt.py} — 在示例2的基础上增加质子两点函数的计算，使用标准的颜色收缩重子算符和宇称投影。

\subsection{示例4：PCAC验证}

\texttt{examples/4\_Pion\_PCAC.py} — 通过计算π介子轴矢量关联函数验证PCAC关系。

\subsection{示例5：π介子色散关系}

\texttt{examples/5\_Pion\_Dispersion.py} — 使用动量投影计算不同动量下的π介子能量，验证连续色散关系$E^2 = m_\pi^2 + p^2$。

%=============================================================================
\section{QUDA底层API}
%=============================================================================

\subsection{参数结构体}

PyQUDA暴露了QUDA的所有主要参数结构体：

\begin{lstlisting}[caption={QUDA参数结构体}]
# 规范场参数
from pyquda import QudaGaugeParam
gauge_param = QudaGaugeParam()
gauge_param.cpu_prec = 2          # 双精度
gauge_param.cuda_prec = 2         # 双精度
gauge_param.reconstruct = 18      # 无重构
gauge_param.t_boundary = -1       # 反周期时间边界

# 反演参数
from pyquda import QudaInvertParam
invert_param = QudaInvertParam()
invert_param.tol = 1e-12          # 收敛精度
invert_param.maxiter = 10000      # 最大迭代次数
invert_param.solve_type = 2       # 直接求解 (Direct)
invert_param.solution_type = 1    # 无预处理 (Mat)
invert_param.dslash_type = 1      # Wilson dslash
invert_param.residual_type = 2    # L2相对残差

# 本征求解器参数
from pyquda import QudaEigParam

# 多重网格参数
from pyquda import QudaMultigridParam
\end{lstlisting}

\subsection{主要QUDA函数}

\begin{lstlisting}[caption={直接调用的QUDA函数}]
from pyquda import (
    # 初始化
    initQuda, endQuda, initCommsGridQuda,

    # 规范场
    loadGaugeQuda, saveGaugeQuda, freeGaugeQuda,
    loadCloverQuda, freeCloverQuda,
    gaussGaugeQuda, gaussMomQuda,
    projectSU3Quda,
    computeGaugeForceQuda,
    updateGaugeFieldQuda,

    # 观测量
    plaqQuda, polyakovLoopQuda, gaugeObservablesQuda,

    # 求解器
    invertQuda, invertMultiSrcQuda, invertMultiShiftQuda,
    dslashQuda, dslashMultiSrcQuda,
    MatQuda, MatDagMatQuda,
    cloverQuda,

    # 多重网格 / 本征求解
    newMultigridQuda, destroyMultigridQuda, updateMultigridQuda,
    eigensolveQuda, newDeflationQuda, destroyDeflationQuda,

    # 涂抹 / 流
    performGaugeSmearQuda, performWFlowQuda,
    performWuppertalnStep, performTwoLinkGaussianSmearNStep,

    # 力计算
    computeCloverForceQuda, computeHISQForceQuda,

    # 规范固定
    computeGaugeFixingOVRQuda, computeGaugeFixingFFTQuda,

    # 其他
    momActionQuda, computeKSLinkQuda,
    contractQuda, blasGEMMQuda,
    flushChronoQuda, setVerbosityQuda,
)
\end{lstlisting}

%=============================================================================
\section{高级功能}
%=============================================================================

\subsection{动量相位}

\begin{lstlisting}[caption={动量相位}]
from pyquda_utils.phase import MomentumPhase, GridPhase

# 动量投影相位
phase = MomentumPhase([24, 24, 24, 72], latt_info, mom=[0, 0, 1, 0])
prop_p = phase @ prop  # 动量投影传播子

# 格点步幅相位
grid_phase = GridPhase(latt_info, step=[2, 2, 2, 1])
\end{lstlisting}

\subsection{Wilson圈计算}

\begin{lstlisting}[caption={Wilson圈}]
from pyquda_utils.wilson_loop import wilson_loop

W = wilson_loop(gauge, R=3, T=4, dirs=["X", "Y", "Z"])
\end{lstlisting}

\subsection{规范场转换}

\begin{lstlisting}[caption={规范场格式转换}]
from pyquda_utils.convert import gauge_convert

# SU(3)规范场转换为SU(2)子群
gauge_su2 = gauge_convert(gauge, "su3_to_su2")
\end{lstlisting}

\subsection{FFT工具}

\begin{lstlisting}[caption={FFT操作}]
from pyquda_utils.fft import fft, ifft

# 格点场的快速傅里叶变换
freq = fft(fermion_field)
\end{lstlisting}

%=============================================================================
\section{API接口清单}
%=============================================================================

\subsection{初始化模块（pyquda/\_\_init\_\_.py）}

\begin{lstlisting}[caption={初始化函数}]
def init(
    resource_path: str,
    grid_type: str = "default",
    device: int = 0,
    enable_quda: bool = True,
    **grid_kwargs,
) -> None: ...

def initGrid(grid_size: list[int], grid_type: str = "default"): ...
def initDevice(device: int): ...
def initQUDA(): ...
def isQUDAInitialized() -> bool: ...
def end(): ...
\end{lstlisting}

\subsection{场模块（pyquda\_comm/field.py）}

\begin{lstlisting}[caption={场基类}]
class LatticeInfo:
    def __init__(self, dims, grid_size=None, t_boundary=-1, anisotropy=1.0): ...
    # 属性: Gs, Gt, volume, grid_size, dims

class LatticeGauge(BaseField):
    def loadNPYGauge(self, path): ...
    def saveNPYGauge(self, path): ...
    def loadHDF5Gauge(self, path): ...
    def saveHDF5Gauge(self, path): ...
    def smearAPE(self, coeff, steps): ...
    def smearStout(self, coeff, steps): ...
    def smearHYP(self, coeff, steps): ...
    def wflow(self, steps, step_size): ...
    def symanzikFlow(self, steps, step_size): ...
    def gaugeFixingOVR(self, tolerance): ...
    def gaugeFixingFFT(self): ...
    def plaquette(self) -> float: ...
    def polyakov(self) -> float: ...
    def topologicalCharge(self) -> float: ...
    def energy(self) -> float: ...
    @property
    def gauge_dirac(self) -> GaugeDirac: ...

class LatticeFermion(BaseField): ...
class LatticePropagator(BaseField): ...
class LatticeMom(BaseField): ...
class LatticeClover(BaseField): ...
class LatticeRotation(BaseField): ...
\end{lstlisting}

\subsection{Dirac算子模块（pyquda/dirac/）}

\begin{lstlisting}[caption={Dirac算子接口}]
class Dirac:
    def __init__(self, latt_info, gauge_param=None, invert_param=None): ...

class FermionDirac(Dirac):
    def invert(self, source, mass=None, tol=1e-12, maxiter=10000): ...
    def invertRestart(self, source, mass=None): ...
    def invertPC(self, source, mass=None): ...
    def invertCloverPC(self, source, mass=None): ...
    def invertMultiSrc(self, sources, mass=None): ...
    def invertMultiSrcRestart(self, sources, mass=None): ...
    def mat(self, vector): ...
    def matDagMat(self, vector): ...
    def dslash(self, vector): ...
    def performance(self): ...

class WilsonDirac(FermionDirac): ...
class CloverWilsonDirac(FermionDirac): ...
class StaggeredDirac(StaggeredFermionDirac): ...
class HISQDirac(StaggeredFermionDirac): ...

class GaugeDirac(Dirac):
    def covDev(self, field, dir, displacement): ...
    def laplacian(self, field): ...
    def smearWuppertal(self, source, n_steps, sigma): ...
    def smearGaussian(self, source, n_steps, sigma): ...
    def projectSU3(self, gauge): ...
    def gaugePath(self, path): ...
    def gaugeLoopTrace(self, path): ...
    def wflow_scale_determination(self, gauge, steps, step_size): ...
\end{lstlisting}

\subsection{作用量模块（pyquda/action/）}

\begin{lstlisting}[caption={作用量接口}]
class Action:
    def action(self) -> float: ...
    def force(self) -> float: ...

class GaugeAction(Action):
    def __init__(self, latt_info, loop_param: LoopParam, beta: float): ...

class FermionAction(Action):
    def __init__(self, latt_info, gauge, mass, n_flavor, ...): ...
    def sample(self): ...
    def invertMultiShift(self, source, masses): ...

class CloverWilsonAction(FermionAction):
    def __init__(self, latt_info, gauge, mass, clover_coeff, n_flavor): ...

class HISQAction(StaggeredFermionAction):
    def __init__(self, latt_info, gauge, mass, n_flavor): ...

class LoopParam:
    def __init__(self, paths: list[list[int]], coeffs: list[float]): ...
\end{lstlisting}

\subsection{HMC模块（pyquda/hmc.py）}

\begin{lstlisting}[caption={HMC接口}]
class HMC:
    def __init__(self, latt_info, monomials: list[Action], integrator: Integrator): ...
    def initialize(self, seed: int, gauge: LatticeGauge): ...
    def gaussMom(self): ...
    def samplePhi(self): ...
    def integrate(self, t: float, tol: float): ...
    def accept(self, dE: float) -> bool: ...
    def gaugeAction(self) -> float: ...
    def fermionAction(self) -> float: ...
    def momAction(self) -> float: ...
    def plaquette(self) -> float: ...
    def polyakovLoop(self) -> float: ...
    def energy(self) -> float: ...
    def qcharge(self) -> float: ...
    def qchargeDensity(self) -> float: ...
    def saveGauge(self, gauge: LatticeGauge): ...
    def loadGauge(self, gauge: LatticeGauge): ...

class Integrator:
    def __init__(self, n_steps: int): ...
    def integrate(self, updateGauge, updateMom, t: float): ...
\end{lstlisting}

\subsection{Gamma矩阵模块（pyquda/gamma.py）}

\begin{lstlisting}[caption={Gamma矩阵接口}]
class Gamma:
    def __init__(self, index: int, factor=1.0): ...
    # index: 0-15 (DeGrand-Rossi基)
    # 支持: +, -, *, @, .T (转置), .H (厄米共轭), .D (dagger)
    @property
    def matrix(self): ...       # 4x4 numpy数组
    @property
    def sparse_indices(self): ...
    @property
    def sparse_data(self): ...

class Projector:
    def __init__(self, factors: list[float]): ...  # 长度16
    # 支持: +, -, *, @, .T, .H
\end{lstlisting}

\subsection{源构造模块（pyquda\_utils/source.py）}

\begin{lstlisting}[caption={源构造接口}]
def point(latt_info, position): ...         # 点源
def wall(dims, latt_info, t=0): ...         # 墙源
def volume(dims, latt_info): ...            # 体积源
def momentum(latt_info, mom): ...           # 动量源
def colorvector(latt_info, color, spin): ... # 颜色矢量源
def source(latt_info): ...                  # 一般源
def source12(latt_info, source_type=None): ... # 12分量源（自旋+颜色）
def source3(latt_info, source_type=None): ...  # 3分量源（交错费米子）
\end{lstlisting}

\subsection{I/O模块（pyquda\_io/）}

\begin{lstlisting}[caption={I/O接口}]
# chroma.py
def readChromaGauge(path: str, latt_info) -> LatticeGauge: ...
def readChromaPropagator(path: str, latt_info) -> LatticePropagator: ...

# milc.py
def readMILCGauge(path: str, latt_info) -> LatticeGauge: ...
def writeMILCGauge(path: str, gauge): ...
def readMILCPropagator(path: str, latt_info) -> LatticePropagator: ...

# npy.py
def readNPYGauge(path: str, latt_info) -> LatticeGauge: ...
def writeNPYGauge(path: str, gauge): ...

# kyu.py, xqcd.py, nersc.py, openqcd.py, ildg.py
# 类似的 read/write 接口

# general.py
def readGeneralGauge(path: str, latt_info) -> LatticeGauge: ...
def writeGeneralGauge(path: str, gauge): ...
\end{lstlisting}

\subsection{工具模块（pyquda\_utils/core.py）}

\begin{lstlisting}[caption={core.py 便捷函数}]
def init(resource_path=".cache/quda", grid_type="default", device=0): ...
def invert(dirac, source, mass=None, tol=1e-12, maxiter=10000): ...
def getDirac(latt_info, gauge, clover_coeff=None, mass=-0.2770): ...
def getWilson(latt_info, gauge, mass=-0.2770): ...
def getClover(latt_info, gauge, mass=-0.2770, clover_coeff=1.0, multigrid=True): ...
def getHISQ(latt_info, gauge, mass=0.05): ...
def getStaggeredDirac(latt_info, gauge, mass=0.05): ...
def gatherLattice(field) -> np.ndarray: ...   # MPI收集+归约
\end{lstlisting}

\subsection{插件模块（pyquda\_plugins/）}

\begin{lstlisting}[caption={插件接口}]
# pycontract
def mesonTwoPoint(prop1, prop2, gamma_list): ...
def mesonAllSinkTwoPoint(prop1, prop2, gamma_list): ...
def baryonTwoPoint(prop1, prop2, gamma_list): ...
def baryonDiquark(prop1, prop2, gamma_list): ...
def baryonSequentialTwoPoint(prop1, prop2, seq_prop, gamma_list): ...

# pygwu
class Overlap:
    def __init__(self, latt_info, gauge, mass, kappa): ...
    def buildWilsonKernel(self): ...
    def invert(self, source, mass, tol=1e-12): ...
\end{lstlisting}

%=============================================================================
\section{性能基准测试}
%=============================================================================

PyQUDA的基准测试说明参见项目Wiki：
\url{https://github.com/CLQCD/PyQUDA/wiki/Benchmark}

QUDA环境变量（30+）可用于精细调节GPU性能：
\begin{itemize}[leftmargin=*]
    \item \texttt{QUDA\_RESOURCE\_PATH}：QUDA资源缓存路径
    \item \texttt{QUDA\_ENABLE\_P2P}：启用GPU对等通信
    \item \texttt{QUDA\_ENABLE\_GDR}：启用GPU直连RDMA
    \item \texttt{QUDA\_MPI\_REDUCE\_LOCAL\_COPY}：MPI归约的本地拷贝优化
\end{itemize}

%=============================================================================
\section{维护说明}
%=============================================================================

\begin{itemize}[leftmargin=*]
    \item 函数定义（主要在\texttt{quda.in.pxd}和\texttt{quda.in.pyx}中）及大部分文档字符串（主要在\texttt{pyquda.pyi}和\texttt{enum\_quda.in.py}中）需要手动更新，因为目前无法自动生成
    \item PyQUDA应与QUDA未来的API兼容版本配合良好，前提是QUDA的API保持不变
    \item 核心Cython绑定（\texttt{quda.pyx}）从QUDA的C头文件使用pycparser自动生成，插件系统也采用同样的机制
\end{itemize}

\subsection{更新QUDA API}

当QUDA更新时，需要手动同步以下文件：
\begin{itemize}[leftmargin=*]
    \item \texttt{pyquda/quda.pyi}：Python类型存根
    \item \texttt{pyquda/quda.in.pxd}：Cython声明
    \item \texttt{pyquda/quda.in.pyx}：Cython实现
    \item \texttt{pyquda/enum\_quda.in.py}：枚举类型
\end{itemize}

\end{document}
